% ==========================================================================
% Supervisor Meeting Agenda - 30 March 2026 (30 minutes)
%
% Supervisors: Dr. Josef Bajada (Academic), Dr. Graham Hili (Industry)
% Compile with: pdflatex agenda.tex
% ==========================================================================

\documentclass[11pt,a4paper]{article}
\usepackage[utf8]{inputenc}
\usepackage[margin=2.5cm]{geometry}
\usepackage{enumitem}
\usepackage{hyperref}
\usepackage{booktabs}
\usepackage{tabularx}
\usepackage{xcolor}

\definecolor{done}{RGB}{34,139,34}
\definecolor{todo}{RGB}{200,50,50}
\definecolor{question}{RGB}{0,90,180}

\newcommand{\done}[1]{\textcolor{done}{\textbf{DONE:}} #1}
\newcommand{\todo}[1]{\textcolor{todo}{\textbf{TODO:}} #1}
\newcommand{\question}[1]{\textcolor{question}{\textbf{Q:}} #1}

\title{Supervisor Meeting Agenda\\
\large 30 March 2026 --- 30 minutes}
\author{Antonio Galdes}
\date{}

\begin{document}
\maketitle

% ------------------------------------------------------------------
\section*{1. Progress Update (5 min)}
% ------------------------------------------------------------------

\begin{itemize}[leftmargin=*]
  \item \done{Training pipeline verified end-to-end: 10\,000-step smoke run completed
    successfully. Evidential metrics confirmed live during training
    (epistemic uncertainty = 1.44, aleatoric = 2.07). All three containers
    healthy throughout.}

  \item \done{Sensor-to-covariance pipeline fully validated. Dry-run inspection
    (8 episodes, 300 steps each) confirmed: EKF pose tracks CARLA ground truth
    to within $\sim$1\,m relative error; odom-frame target bay transform is
    geometrically correct (round-trip reconstruction error = 0.000\,m);
    covariance inflation events (staleness, Cartographer jump) visible in log
    output.}

  \item \done{Critical pipeline bugs resolved since last meeting:
    \begin{itemize}
      \item \texttt{imu0\_config} yaw absolute flag (index 5) was incorrectly
        set to \texttt{true} --- EKF was treating IMU as an absolute heading
        source. Fixed to \texttt{false}; only yaw rate (index 11) enabled.
      \item \texttt{twist\_in\_odom\_frame} parameter was defined in
        \texttt{ros2\_config.yaml} but not wired into the launch file.
        Now propagated correctly to \texttt{CovarianceExtractorNode}.
      \item Odom-frame target bay transform implemented
        (\texttt{\_compute\_target\_bay\_odom}): both vehicle pose and target
        are expressed in the same accumulated Cartographer odom frame, so
        mid-episode drift cancels in the relative target vector
        (obs[15--17]).
    \end{itemize}}

  \item \done{Placement work synthesised. The placement objective (preliminary
    design of an uncertainty-aware RL policy) has been exceeded --- a full
    working implementation was produced. Placement demo and one-page report
    are being prepared (see Section~2.1).}

  \item \todo{Full 1M-step training run not yet launched --- pipeline is
    verified and ready. Covariance variation under NPC vehicles and weather
    randomisation still to be confirmed during live training.}

  \item \todo{Mapping runs for \texttt{trapezoid} and \texttt{irregular\_a}
    floor plans still outstanding ($\sim$90\,s each,
    \texttt{make docker-map LAYOUT=<name>}).}
\end{itemize}

% ------------------------------------------------------------------
\section*{2. Discussion Points (15 min)}
% ------------------------------------------------------------------

\subsection*{2.1 Placement Deliverables --- Demo and One-Page Report}

The placement brief (Lemonworx, 8 weeks) required a preliminary design for an
uncertainty-aware RL policy. The implementation is substantially complete. Two
placement deliverables remain:

\textbf{Demo plan (approx.\ 10 minutes, live):}
\begin{enumerate}
  \item \texttt{make docker-up} --- show the three-container stack is healthy.
  \item \texttt{make docker-inspect INSPECT\_LAYOUT=trapezoid} --- CARLA window
    with parking lot geometry, patrol NPC, and pedestrians.
  \item \texttt{make docker-inspect-sensors INSPECT\_SUITE=suite\_a} --- LiDAR
    arc and IMU position on the ego vehicle, connecting sensor to uncertainty signal.
  \item Code walkthrough: \texttt{evidential\_policy.py} (NIG output head,
    epistemic/aleatoric decomposition) and \texttt{sb3\_integration.py}
    (evidential regularisation term in PPO loss).
  \item Live ROS\,2 topic echo of \texttt{/ekf\_uncertainty/covariance} showing
    uncertainty values changing as simulation conditions change.
\end{enumerate}

\textbf{One-page report structure:}
(1) literature gap identified; (2) preliminary architecture design produced
(actor-only EDL + EKF covariance as explicit state input); (3) implementation
delivered beyond the placement scope; (4) relevance to teleoperation cognitive
load problem (uncertainty thresholds trigger operator handoff).

\question{Is this demo scope appropriate for the placement assessment, or should
it focus more on the theoretical design decisions rather than the live system?}

\question{The one-page report will reference the 2$\times$2 ablation as the
evaluation strategy. Should it also include a brief statement on the
sim-to-real transfer plan (same Docker stack, physical BMW\,i3 sensors), or
is that scope beyond the placement brief?}

\bigskip
\subsection*{2.2 Covariance Variation --- Sufficiency for Training}

Dry-run inspection with random actions in a clear empty lot shows baseline
std values of $\sim$0.10--0.12\,m (flat, as expected with no occlusion).
Meaningful covariance variation requires NPC vehicles and weather randomisation
active during a real training run. This has not yet been observed live.

\question{Is it acceptable to begin the full 1M-step training run before
separately verifying covariance variation, or should a short dedicated
verification run (NPC vehicles + HardRainNoon, monitor covariance only,
no training) be conducted first to confirm the signal is present?}

\question{If covariance variation proves insufficient under the current
sensor suite (Suite A: 2D LiDAR only), what is the preferred fallback:
(a) Suite B (16-channel 3D LiDAR, wider scan coverage), (b) artificial
noise injection on the EKF input, or (c) additional NPC density beyond the
current maximum of 3 patrol vehicles and 4 pedestrians?}

\bigskip
\subsection*{2.3 Cartographer Odom Frame Accumulation Across Episodes}

The Cartographer odom frame accumulates continuously across training episodes
(it is a process-level frame, not per-episode). As a result, absolute odom
coordinates grow over many episodes. This is not a bug --- the relative target
vector (obs[15--17]) remains correct because drift cancels --- but it means
the raw pose values in obs[0--2] grow without bound.

\texttt{VecNormalize} handles this via running statistics. The policy never
sees unbounded values.

\question{Is it acceptable to let absolute odom coordinates grow unbounded
(mitigated by VecNormalize), or is a per-episode odom frame reset needed?
Resetting would require restarting the Cartographer ROS process per episode
($\sim$10--30\,s overhead), which is likely impractical for 1M-step training.}

% ------------------------------------------------------------------
\section*{3. Timeline and Next Steps (5 min)}
% ------------------------------------------------------------------

\begin{table}[h]
\centering
\small
\begin{tabular}{@{}lll@{}}
\toprule
\textbf{Task} & \textbf{Est.\ Time} & \textbf{Status} \\
\midrule
Mapping runs (trapezoid, irregular\_a) & 0.5 days & Ready to run \\
Full 1M-step training run & 2--4 days (GPU) & Unblocked \\
Covariance variation verification & During training & Passive monitoring \\
Reward tuning (Task 9) & 1--2 weeks & After first run \\
Full ablation (4 baselines $\times$ 10 seeds) & 2--4 weeks & After Task 9 \\
Tasks 10--14: Evaluation + plots & 2 weeks & After ablation \\
Placement demo + one-page report & 2--3 days & In progress \\
Dissertation writing & Ongoing & Can start now \\
\bottomrule
\end{tabular}
\end{table}

\question{Does this timeline remain on track for the September 2026 submission?
Are there any tasks that should be deprioritised or cut if the ablation
study takes longer than expected?}

% ------------------------------------------------------------------
\section*{4. Dissertation Writing (5 min)}
% ------------------------------------------------------------------

\begin{itemize}[leftmargin=*]
  \item The methodology chapter is ready to draft: architecture is stable,
    all design decisions are resolved, and the sensor-to-covariance pipeline
    is fully documented.
  \item The background chapter (EKF, evidential deep learning, PPO,
    Cartographer) is independent of experimental results and can be written
    in parallel with training.
  \item The experimental setup chapter is also ready: floor plans, sensor
    suite, NPC configuration, and the Cartographer pure localisation workflow
    are all fixed.
\end{itemize}

\question{Any guidance on chapter structure, expected length, and when to
submit draft chapters for feedback? Should chapters be submitted individually
as they are written, or as a complete draft?}

\question{The placement proposal and dissertation proposal share the same
research problem and evaluation strategy. Should the dissertation introduction
explicitly acknowledge the placement as the origin of the implementation, or
treat the two submissions as entirely separate documents?}

\end{document}
